%%%%%%%%%%%%%%%%%%%%%% --- Informacje dla studentów --- %%%%%%%%%%%%%%%%%%%%%%%
%--- Informacje rozpoczynające się od znaków "&---" są informacjami dla studenta dotyczącymi edycji tekstu.

%--- Część ustawień dokumentu została wprowadzona jako informacja alternatywna &\. Funkcje w tekście zamiast zaczynać się od backslasha "\" zaczyna się od "%\". W takiej sytuacji należy usunąć %, aby uruchomić nową opcje dokumentu, a w starej wstawić %.

%--- Przed rozpoczęciem edycji pliku w Overleaf należy wejść w Menu > Spell check > Polish, aby włączyć sprawdzanie tekstu w języku polskim.

%--- Akapity wypełnione tekstem "Lorem ipsum dolor sit amet, consectetue..." symulują przepływ tekstu w dokumencie i powinny zostać usunięte lub zastąpione tekstem pracy dyplomowej.

%--- Wzór pracy dyplomowej wykorzystuje bibliografię automatyczną z pliku BibTeX. Przykładowa bibliografia znajduje się z pliku bib.bib. Plik ten można edytować lub zamienić pod tą samą nazwą. Do utworzenia tego pliku najlepiej użyć menedżera bibliografii Mendeley lub Zotero. Szczegóły odnośnie formatowania bibliografii i odwołań w tekście znajdują się we wzorze.

%--- Wzór pracy dyplomowej przygotował Marek Deja z Instytutu Studiów Informacyjnych UJ, kontakt: marek.deja@uj.edu.pl

%%%%%%%%%%%%%%%%%%%%%%% --- Ustawienia dokumentu. --- %%%%%%%%%%%%%%%%%%%%%%%
\documentclass[twoside,a4paper,12pt]{report} %--- Druk dwustronny: twoside. Druk jednostronny: ondeside
\usepackage{polski}
\usepackage[utf8]{inputenc}
\usepackage{indentfirst}
\usepackage{microtype}
\DisableLigatures{encoding = *, family = * }
\usepackage{fancyhdr}
\usepackage{amssymb}
\usepackage{url,hyperref}
\usepackage{multirow,multicol}
\usepackage[singlelinecheck=false]{caption}
\usepackage{booktabs,float,tabularx,ltablex}
\usepackage[flushleft]{threeparttable}
\usepackage{enumitem}
\usepackage{pstricks,graphicx}
\usepackage{amsmath}
\usepackage[nottoc]{tocbibind}
\usepackage{natbib}
\bibliographystyle{apalike}%--- Styl bibliograficzny APA 6th
\usepackage{lipsum}
\usepackage{chngcntr}
\counterwithout{table}{chapter}
\counterwithout{figure}{chapter}
\usepackage{titlesec}
\usepackage{tikz}
\usepackage{pgfplots}
\pgfplotsset{compat=1.14}
\floatstyle{plaintop}
\newfloat{wykres}{tbph}{loc}
\floatname{wykres}{Wykres}
\usepackage{cleveref}
\crefname{wykres}{wykres}{wykres}
\Crefname{wykres}{wykres}{Wykres}
\newcommand\listofcharts{\listof{wykres}{Spis wykresów}}
\usepackage{geometry}
\newgeometry{tmargin=2.5cm, bmargin=2.5cm, headheight=16pt, inner=3.5cm, outer=2cm} 
\usepackage{pdflscape}
\usepackage{rotating}
\usepackage{theorem}
\theoremstyle{break}
\theorembodyfont{\it}
\newtheorem{hip}{Hipoteza}[chapter]
\theorembodyfont{\rm}
\newtheorem{defi}{Definicja}[chapter]

%%%%%%%%%%%%%%%%%%%%%%% --- Ustawienia układu strony --- %%%%%%%%%%%%%%%%%%%%%%%

\pagestyle{fancy}
%---Druk jednostronny. Aby wyłączyć należy dodać % przed \. Aby włączyć należy usunąć % prze \.

%\fancyhead[C]{} 
%\fancyfoot[R]{\thepage}
%\fancyhead[L]{\scriptsize\leftmark}
%\fancyhead[R]{\scriptsize\rightmark}
%\cfoot[]{}
%\fancypagestyle{plain}{%
%\fancyhf{} % clear all header and footer fields
%\fancyfoot[R]{\textbf{\thepage}} % except the center
%\renewcommand{\headrulewidth}{0pt}
%\renewcommand{\footrulewidth}{0pt}}

\setcounter{secnumdepth}{4}
%---Druk dwustronny. Aby wyłączyć należy dodać % przed \. Aby włączyć należy usunąć % prze \.

\fancyhead[LE,RO]{\scriptsize\leftmark}
\fancyhead[RE,LO]{{\scriptsize \rightmark}}
\fancyfoot[CE,CO]{}
\fancyfoot[LE,RO]{\thepage}
\cfoot[]{}
\fancypagestyle{plain}{%
\fancyhf{} % clear all header and footer fields
\fancyfoot[LE,RO]{\textbf{\thepage}} % except the center
\renewcommand{\headrulewidth}{0pt}
\renewcommand{\footrulewidth}{0pt}}

%%%%%%%%%%%%%%%%%%%%%%% --- Koniec edycji układu strony. Proszę nie edytować --- %%%%%%%%%%%%%%%%%%%%%%%
\renewcommand{\chaptermark}[1]{%
\markboth{\MakeUppercase{%
\chaptername}\ \thechapter.%
\ #1}{}}
\renewcommand{\sectionmark}[1]{\markright{\thesection.\ #1}}

%%%%%%%%%%%%%%%%%%%%%% --- Wypełnia student --- %%%%%%%%%%%%%%%%%%%%%%%

\title{Tytuł rozprawy}
\author{Imię i Nazwisko }
\date{Maj 2020}


%%%%%%%%%%%%%%%%%%%%%%% --- Formatowanie strony tytułowej --- %%%%%%%%%%%%%%%%%%%%%%%
\begin{document}
\titleformat{\chapter}[display] {\normalfont\Huge\bfseries}{\chaptertitlename\ \thechapter}{0pt}{\Huge}
\titlespacing{\chapter}{1cm}{2cm}{1cm}

\thispagestyle{empty}
\begin{center}{\sc \large
Uniwersytet Jagielloński\par\vspace{0.2cm}\par
Wydział Zarządzania i Komunikacji Społecznej}\par\vspace{0.3cm}\par
{\LARGE
INSTYTUT STUDIÓW INFORMACYJNYCH}\par\vspace{0.2cm}\par
{\large
Studia stacjonarne}\par\vspace{1cm}\par
{\Large
Praca magisterska}%--- Wypełnia student --- "Praca licencjacka" w przypadku prac licencjackich 
\end{center}
\vspace{1.5cm}
\begin{flushleft}
{\large
Nr albumu: XXXXXX}%--- Wypełnia student
\end{flushleft}
\vspace{1.5cm}
\begin{center}
{\LARGE
Jan Kowalski%--- Wypełnia student
}\par\vspace{0.9cm}\par
{\huge
TYTUŁ PRACY%--- Wypełnia student
}
\end{center}
\vspace{4cm}
\begin{flushright}

Opiekun pracy magisterskiej:\\
dr hab. Piotr Kowalski%--- Wypełnia student
\end{flushright}
\vfill
\begin{center}
Kraków 2020%--- Wypełnia student
\end{center}


\leavevmode\thispagestyle{empty}\newpage
\centerline{\LARGE \bf Streszczenie}\par\vspace{1cm}\par
\lipsum[1]
\thispagestyle{empty} 
\newpage
\centerline{\LARGE \bf Summary}\par\vspace{1cm}\par
\lipsum[1]
\thispagestyle{empty} 
\newpage

\vspace*{16cm}
\begin{flushright}
miejsce na ewentualne podziękowania\\
dla promotora%-- -Wypełnia studen. Można usunąć ten tekst i pozostawić pustą stronę
\end{flushright}

\thispagestyle{empty} 
\tableofcontents %-- automatyczny spis treści

%%%%%%%%%%%%%%%%%%%%% --- Zasadnicza treść rozprawy zaczyna się tutaj --- %%%%%%%%%%%%%%%%%%%%%

\chapter*{Wstęp}

\addcontentsline{toc}{chapter}{Wstęp}
Rozprawa powinna zostać przygotowana zgodnie instrukcją zamieszczoną na stronie Instytutu Studiów Informacyjnych.\medskip \\
\url{https://isi.uj.edu.pl/studia/instrukcje}\bigskip

\lipsum[1]%---Zastąpić treścią


\chapter{Nazwa rozdziału}

\lipsum[1-2]\par%---Zastąpić treścią

\section{Nazwa podrozdziału}
\lipsum[1]%---Zastąpić treścią


\subsection{Nazwa podrozdziału}

\lipsum[1]%---Zastąpić treścią

\subsubsection{Nazwa podrozdziału}

\lipsum[1]%---Zastąpić treścią

\begin{hip}[Hipoteza I]

\label{t1}%etykieta hipotezy

Treść hipotezy

\end{hip}

\chapter{Formatowanie tekstu i odwołań}

Dosłowne przywołanie fragmentu innej publikacji lub prac niepublikowanych należy oznaczać w tekście w następujący sposób: cytowany fragment należy umieścić w cudzysłowie, bezpośrednio po którym należy umieścić przypis wskazujący dokładnie źródło cytowania (koniecznie z podaniem numeru strony w przypadku, gdy źródło jest paginowane).\\
Przykład: \par
\sloppy%---komenda używana przed akapitami zawierającymi kursywę, gdy pojawia się problem z łamaniem tekstu.
W {\it Słowniku encyklopedycznym informacji, języków i systemów  informacyjno-wyszukiwawczych} odbiorca komunikatu został zdefiniowany jako: „element układu informacyjnego (uczestnik komunikacji) będący wykonawcą procesu odbierania komunikatu (…)” (Bojar oprac., 2002, s. 177).

\section{Przypisy w tekście i przypisy dolne}

Dokument wykorzystuje {\bf bibliografię automatyczną}, która jest zapisana w pliku \underline{bib.bib}. Można go edytować lub zamienić plikiem o tej samej nazwie, np. wygenerowanym z menedżera bibliografii Mendeley\footnote{\url{http://pka.bj.uj.edu.pl/mendeley/}} lub Zotero\footnote{\url{https://www.zotero.org/}}.\par \bigskip

Przykładowe odwołania w tekście w stylu APA 6th\\
Nazwisko - Rok \citep{Cheuk1998}\\
Nazwisko - Rok - Strona \citep[s. 77]{Collins2010}\\
Rok publikacji \citeyearpar{Choo2013}\\
Wiele odwołań \citep{Ginman1987,Oliver2008}\\
Wiele odwołań ze stronami (\citeauthor{Deja2017}, \citeyear[s. 5]{Deja2017}; \citeauthor{Choo2013}, \citeyear[s. 766]{Choo2013})\par \bigskip

Odwołania do tabel lub wykresów formatujemy w tekście w nawiasach okrągłych, np.: (Tab. ~\ref{tab:tab1}), (Rys. \ref{fig:fig1})


\newpage
\subsection{Listy}
Przykład listy z punktorami:
\begin{itemize}
    \item Pozycja pierwsza.
    \item Pozycja druga.
    \item Pozycja trzecia
    \begin{itemize}
        \item Podpunkt pierwszy.
        \item Podpunkt drugi.
    \end{itemize}
\end{itemize}\bigskip

Przykład listy numerowanej:
\begin{enumerate}
    \item Pozycja pierwsza.
    \item Pozycja druga.
    \item Pozycja trzecia
    \begin{enumerate}
        \item Podpunkt pierwszy.
        \item Podpunkt drugi.
    \end{enumerate}
\end{enumerate}\bigskip

Przykład listy numerowanej:
\begin{enumerate}[label=\roman*.]
    \item Pozycja pierwsza.
    \item Pozycja druga.
    \item Pozycja trzecia
\end{enumerate}\bigskip

Przykład listy numerowanej:
\begin{enumerate}[label=\Roman*.,start=4]
    \item Pozycja pierwsza.
    \item Pozycja druga.
    \item Pozycja trzecia
\end{enumerate}\bigskip

Przykład listy numerowanej:
\begin{enumerate}[label=\alph*.]
    \item Pozycja pierwsza.
    \item Pozycja druga.
    \item Pozycja trzecia
\end{enumerate}\bigskip

Przykład listy numerowanej:
\begin{enumerate}[label=\Alph*.,start=4]
    \item Pozycja pierwsza.
    \item Pozycja druga.
    \item Pozycja trzecia
\end{enumerate}\bigskip

\subsubsection{Dodatkowe zasady}
Tytuły czasopism, książek, zasobów WWW występujące w tekście należy pisać {\it kursywą} bez cudzysłowu (uwaga: inne zasady dotyczą opisów bibliograficznych w przypisach i bibliografii załącznikowej).
\par
Wyrazy w językach obcych należy podawać kursywą. Przykład: {\it novum} 
\par
Wyróżnienia w tekście należy zaznaczać wytłuszczeniem. Przykład: {\bf królowa} \par

\newpage
\section{Tabele, wykresy i rysunki}

Osoby przyzwyczajone do edycji tabel w Excelu mogą posłużyć się dodatkiem do Excela Exce2LaTeX, który można pobrać pod tym adresem: \url{https://ctan.org/tex-archive/support/excel2latex?lang=en}\bigskip

\lipsum[1]%---Zastąpić treścią

\begin{table}[ht]%---w nawiasie kwadratowym parametry pozycji tabeli w tekście
\caption{Przykładowa tabela}
\label{tab:tab1}
%\centering %---wyśrodkowanie tabeli
\begin{tabular}{|l|l|p{5cm}|}%--- Trzy parametry: l - kolumna do wyjustowana lewej, r - kolumna do prawej, c - wyśrodkowana, p - zawijanie do szerokości
    \hline
    Nagłówek1 & Nagłówek2 & Nagłówek3 \\ \hline
    Wiersz1 & Dane1 & Dane1 \\ \hline
    Wiersz1 & Dane2 & Dane2  \\ \hline
    Wiersz1 & Dane3 & Dane3 \\ \hline
\end{tabular}
\begin{tablenotes}
      \small
      \item Źródło: oprac. własne
    \end{tablenotes}

\end{table}

\lipsum[1]

\begin{table}[ht]
  \begin{threeparttable}
    \caption{Tabela dzielona na trzy}
     \label{tab:tab2}
     \begin{tabular}{lllll}
        \toprule
        Nagłówek & \( df \) & \( f \) & \( \eta \) & \( p \) \\
        \midrule
                 &  \multicolumn{4}{c}{Dane}     \\
        Wiersz1    & 1        & 0.67    & 0.55       & 0.41    \\
        Wiersz2    & 2        & 0.02    & 0.01       & 0.39    \\
        Wiersz3    & 3        & 0.15    & 0.33       & 0.34    \\
        Wiersz4    & 4        & 1.00    & 0.76       & 0.54    \\
        \bottomrule
     \end{tabular}
    \begin{tablenotes}
      \small
      \item Źródło: oprac. własne
    \end{tablenotes}
  \end{threeparttable}
\end{table}

\lipsum[1]

\begin{table*}[ht]
\caption{Tabela ze złączonymi wierszami}
\label{tab:tab3}
\begin{tabular}{|l|l|p{4cm}|}
\hline
    Nagłówek1 & Nagłówek2 & Nagłówek3 \\ \hline
   \multirow{2}{*}{Wiersz1} & Dane1 & Dane1 \\
   & Dane2 & Dane2  \\ \hline
    Wiersz2 & Dane3 & Dane3 \\ \hline
\end{tabular}
\begin{tablenotes}
      \small
      \item Źródło: oprac. własne
    \end{tablenotes}
\end{table*}

\lipsum[1]

\section{Rysunki ilustracje}

\begin{figure}[ht]
    \centering
    \caption{Przykładowy rysunek}\label{fig:fig1}
        \includegraphics[width=0.8\textwidth]{fig1.png}
    \begin{tablenotes}
          \item Źródło: oprac. własne
    \end{tablenotes}
\end{figure}

\newpage

\subsection{Wykresy}

Wykresy można dodawać na dwa sposoby. Proste wykresy słupkowe można dodać jak na poniższym przykładzie (Wyk. \ref{wyk1}).\par
Znacznie prościej jest jednak używać wykresów z Excela zapisanych jako PDF (Wyk.\ref{wyk2}) \par
W Excelu należy zaznaczyć wykres, następnie  wybrać "Zapisz jako", PDF, Opcje,"Zaznaczony wykres" i Zapisać. Następnie w Overleaf Wybrać Upload i użyć PDF jak w przypadku wykresu drugiego (poniżej).

\begin{wykres}
\caption{Przykładowy wykres}\label{wyk1}
\begin{tikzpicture}
\begin{axis}[
	x tick label style={
		/pgf/number format/1000 sep=},
	ylabel=Rok,
	enlargelimits=0.05,
	legend style={at={(0.5,-0.2)},
	anchor=north,legend columns=-1},
	ybar interval=0.7,
]
\addplot 
	coordinates {(2012,50) (2011,50)
		 (2010,50) (2009,50)};
\addplot 
	coordinates {(2012,30) (2011,20) 
		(2010,25) (2009,30)};
\legend{Kobiety,Mężczyźni}
\end{axis}
\end{tikzpicture}
\begin{tablenotes}
          \item Źródło: oprac. własne
    \end{tablenotes}
\end{wykres}

\begin{wykres}
\caption{Przykładowy wykres z Excela}\label{wyk2}
    \centering
            \includegraphics[width=0.7\textwidth]{wyk1.pdf}
    \begin{tablenotes}
          \item Źródło: oprac. własne
    \end{tablenotes}
\end{wykres}

\chapter*{Zakończenie}

\addcontentsline{toc}{chapter}{Zakończenie}


%%%%%%%%%%%%%%%%%%%%%% --- Bibliografia i indeksy --- %%%%%%%%%%%%%%%%%%%%%%

\bibliography{bib}%--- Jeżeli okażesię to konieczne, to spolszczenia bibliografii można dokonać darmowym edytorem PDF sejda.com - używając opcji Text i Without, aby zamienić In na W lub editor. na red.


\listoftables
\listoffigures
\listofcharts{}

%%%%%%%%%%%%%%%%%%%%%% --- Aneks --- %%%%%%%%%%%%%%%%%%%%%%

\chapter*{Aneks}
\addcontentsline{toc}{chapter}{Aneks}
\lipsum[2]

\begin{landscape}
\begin{table}[ht]
  \begin{threeparttable}
    %\caption{Tabela dzielona na trzy 2}
     %\label{tab:tab4}
     \begin{tabular}{lllllllllll}
        \toprule
        Nagłówek & \( A \) & \( B \) & \( D \) & \( E \)  & \( D \) & \( E \) & \( D \) & \( E \) & \( D \) & \( E \)\\
        \midrule
                 &  \multicolumn{10}{c}{Dane}     \\
        Wiersz1    & 1        & 0.67    & 0.55       & 0.41       & 0.55       & 0.41       & 0.55       & 0.41       & 0.55       & 0.41    \\
        Wiersz2    & 2        & 0.02    & 0.01       & 0.39        & 0.01       & 0.39        & 0.01       & 0.39        & 0.01       & 0.39    \\
        Wiersz3    & 3        & 0.15    & 0.33       & 0.34        & 0.33       & 0.34        & 0.33       & 0.34        & 0.33       & 0.34    \\
        Wiersz4    & 4        & 1.00    & 0.76       & 0.54        & 0.76       & 0.54        & 0.76       & 0.54        & 0.76       & 0.54    \\
        \bottomrule
     \end{tabular}
    \begin{tablenotes}
      \small
      \item Źródło: oprac. własne
    \end{tablenotes}
  \end{threeparttable}
\end{table}
\end{landscape}

\newpage
\begin{table}[ht]
   \begin{turn}{90}
  \begin{threeparttable}
    %\caption{Tabela dzielona na trzy 3}
     %\label{tab:tab5}
     \begin{tabular}{lllllllllll}
        \toprule
        Nagłówek & \( A \) & \( B \) & \( D \) & \( E \)  & \( D \) & \( E \) & \( D \) & \( E \) & \( D \) & \( E \)\\
        \midrule
                 &  \multicolumn{10}{c}{Dane}     \\
        Wiersz1    & 1        & 0.67    & 0.55       & 0.41       & 0.55       & 0.41       & 0.55       & 0.41       & 0.55       & 0.41    \\
        Wiersz2    & 2        & 0.02    & 0.01       & 0.39        & 0.01       & 0.39        & 0.01       & 0.39        & 0.01       & 0.39    \\
        Wiersz3    & 3        & 0.15    & 0.33       & 0.34        & 0.33       & 0.34        & 0.33       & 0.34        & 0.33       & 0.34    \\
        Wiersz4    & 4        & 1.00    & 0.76       & 0.54        & 0.76       & 0.54        & 0.76       & 0.54        & 0.76       & 0.54    \\
        \bottomrule
     \end{tabular}
    \begin{tablenotes}
      \small
      \item Źródło: oprac. własne
    \end{tablenotes}
  \end{threeparttable}
  \end{turn}
\end{table}
%%%%%%%%%%%%%%%%%%%%%% --- Oświadczenie --- %%%%%%%%%%%%%%%%%%%%%%
\newpage\thispagestyle{empty}
\newpage\thispagestyle{empty}
Wzór oświadczenia: \hfill Kraków, dnia wpisać date\par\vspace{1cm}\par
\centerline{{\Large \bf O Ś W I A D C Z E N I E}}
\par\vspace{1.4cm}\par
Ja niżej podpisany/a
\begin{center}
{\sc \bf IMIĘ I NAZWISKO}
\end{center}
\begin{center}
{\large student/ka Instytutu Studiów Informacyjnych\\
Wydział Zarządzania i Komunikacji Społecznej\\
Uniwersytetu Jagiellońskiego w Krakowie
}\par\vspace{0.6cm}\par
oświadczam,\\

że przedkładaną pracę magisterską pt.

\end{center}
\par
\begin{center}
{\sc \sl \large \bf Tytuł pracy magisterskiej}
\end{center}

\begin{center}
    {\large przygotowaną pod opieką dr hab. Piotr Kowalski}\par\vspace{0.4cm}\par
    {\bf napisałem samodzielnie.}
\end{center}
\par\vspace{0.6cm}\par

Oznacza to, że przy pisaniu pracy, poza niezbędnymi konsultacjami, nie korzystałem/am z pomocy innych osób, a w szczególności nie zlecałem/am opracowania rozprawy lub jej istotnych części innym osobom, ani też nie odpisywałem/am tej rozprawy lub jej istotnych części od innych osób.
\par
Jednocześnie przyjmuję do wiadomości, że gdyby powyższe oświadczenie okazało się nieprawdziwe, decyzja o wydaniu mi dyplomu zostanie cofnięta.
\par\vspace{2cm}\par
\begin{flushright}
....................................... \ \ \ \ \ \ \\
{\scriptsize (podpis)\hspace{2.3cm}\ }
\end{flushright}

\end{document}
